\section{Balancing and the constraint identities}\label{sec:geometry}

The round model suggests which quantities to preserve, but the given
end need not have round leaves or a radial mass aspect. We begin by
choosing coordinates in which the Bondi momentum vanishes. This change
does not alter the data; it removes the first spherical harmonic of the
mass aspect, which is the obstruction to the angular divergence equation
used later. We then write the constraints for the original, generally
nonround leaves. These identities will let us round them without losing
the dominant energy condition.

All remainder estimates below have the componentwise, differentiated
meaning in \eqref{eq:cks}. Only three derivatives of the metric
remainder and two of the tensor remainder are used.

\subsection{A balanced end chart}

\begin{lemma}[Lorentz balancing]\label{lem:balance}
Let smooth data $(g,K)$ have an end in the asymptotic class
\eqref{eq:cks}, with charges defined by \eqref{eq:bondi}.
This class is preserved by the coordinate changes
induced by proper orthochronous Lorentz transformations of the background
unit hyperboloid. In the pullback convention described below, the charge
vector transforms by
\[
 (E_B',P_B')=\mathscr L^{-1}(E_B,P_B).
\]
Consequently, if $E_B>|P_B|$, there is an end chart in the same class in
which $E_B'=m_B$ and $P_B'=0$. In this chart the smooth function
\begin{equation}\label{eq:balanced-aspect}
 f_*:=\frac{M}{4}
 \qquad\text{satisfies}\qquad
 \avg{f_*}=m_B,
 \qquad
 \int_{\mathbb S^2}f_*n^i\,\mathrm dA_\sigma=0
 \quad (i=1,2,3).
\end{equation}
Large spheres in the new chart still bound compact inner portions of the
same exterior.
\end{lemma}

\begin{proof}
Use the Minkowski inner product of signature $(-,+,+,+)$ and write
\[
 X=(\sqrt{1+r^2},rn),\qquad
 X'=(\sqrt{1+\rho^2},\rho n'),\qquad X=\mathscr L X'.
\]
Since $\mathscr L$ preserves the future null cone, there are a positive
smooth function $D$ and a sphere diffeomorphism $\Phi$ such that
\begin{equation}\label{eq:lorentz-boundary}
 \mathscr L(1,n')=D(n')(1,\Phi(n')).
\end{equation}
For tangent vectors $Y,Z$ on the primed sphere, the Minkowski
inner product of the derivatives of the left side is $\sigma(Y,Z)$.
On the right side it is $D^2(\Phi^*\sigma)(Y,Z)$: the derivatives of
$D$ multiply a null vector orthogonal to all angular derivatives.
Consequently
\begin{equation}\label{eq:lorentz-conformal}
 \Phi^*\sigma=D^{-2}\sigma,
 \qquad \Phi^*\mathrm dA_\sigma=D^{-2}\mathrm dA_\sigma.
\end{equation}
The area formula follows because the sphere is two-dimensional.
Expanding $X'$ in powers of $\rho^{-1}$ gives
\[
 r=D\rho+O(\rho^{-1}),\qquad
 n=\Phi(n')+O(\rho^{-2}),
\]
with the same orders after any fixed number of angular derivatives and
applications of $\rho\partial_\rho$. In particular,
\[
 \partial_\rho r=D+O(\rho^{-2}),\qquad
 \partial_A r=O(\rho),\qquad
 \partial_\rho n=O(\rho^{-3}),\qquad
 \partial_A n=\partial_A\Phi+O(\rho^{-2}).
\]
The metric $h$ keeps its standard form because it is the metric induced
on the unit hyperboloid.

We verify the perturbation orders directly; this avoids requiring a
stronger expansion than \eqref{eq:cks}. For $a=g-h$, the leading new
radial component is
\[
 a_{\rho\rho}
 =\frac{D^{-3}(m^r\circ\Phi)}{\rho^5}+O_3(\rho^{-6}).
\]
Indeed, the radial--radial term has this coefficient; the old mixed
components contribute $O_3(\rho^{-6})$, and the old angular components
contribute $O_3(\rho^{-7})$. In a new angular--angular component, the
old radial--radial and mixed terms have orders $O_3(\rho^{-3})$ and
$O_3(\rho^{-2})$, respectively. Thus
\[
 a_{AB}
 =\frac{D^{-1}(\Phi^*m^g)_{AB}}{\rho}+O_3(\rho^{-2}),
 \qquad a_{\rho A}=O_3(\rho^{-3}).
\]
The same calculation for $b=K-h$ gives
\[
 b_{\rho\rho}=O_2(\rho^{-5}),\qquad
 b_{\rho A}=O_2(\rho^{-3}),\qquad
 b_{AB}=\frac{D^{-1}(\Phi^*m^K)_{AB}}{\rho}
       +O_2(\rho^{-2}).
\]
These statements include their indicated derivatives. To see this
without losing a derivative, observe that a new scaled radial or
angular derivative is a combination, with uniformly bounded smooth
coefficients, of $r\partial_r$ and old angular derivatives. Apply the
chain and product rules through order three for $a$ and order two for
$b$. The coordinate coefficients themselves have expansions to every
order. No derivative of the original perturbations beyond these orders
is used. Taylor expansion of a leading angular coefficient is harmless:
these coefficients are smooth, even though only the stated finite
number of weighted derivatives of the remainders is assumed.

It follows from \eqref{eq:lorentz-conformal} that the mass aspect in
\eqref{eq:bondi} transforms as
\[
 M'=D^{-3}(M\circ\Phi).
\]
For completeness, the trace in this calculation satisfies
$\tr_\sigma(\Phi^*T)=D^{-2}(\tr_\sigma T)\circ\Phi$ for any
covariant two-tensor $T$ on the sphere. Combining the charge definition
with \eqref{eq:lorentz-boundary} and \eqref{eq:lorentz-conformal}, we obtain
\begin{align*}
 \frac{1}{16\pi}\int_{\mathbb S^2}(1,n')M'\,\mathrm dA_\sigma
 &=\frac{1}{16\pi}\int_{\mathbb S^2}
       D^{-1}(1,n')(M\circ\Phi)\,\Phi^*\mathrm dA_\sigma\\
 &=\mathscr L^{-1}\frac{1}{16\pi}\int_{\mathbb S^2}
       (1,n)M\,\mathrm dA_\sigma.
\end{align*}
This proves the transformation law for the full charge, including its
extrinsic-curvature contribution. It agrees with the charge conventions
of \cite[Remark~2.1]{CKS} and \cite[Definitions~1.3--1.4]{CWY}.

Strict timelikeness permits the choice
$\mathscr L(m_B,0)=(E_B,P_B)$, which proves
\eqref{eq:balanced-aspect}. Since $D$ is positive on the compact sphere,
$r$ and $\rho$ are uniformly comparable sufficiently far out. The new
chart therefore contains a full tail, and its sufficiently large
spheres have compact inner sides, just as in the original end chart.
\end{proof}

We use this balanced chart from now on and again call its coordinates
$(r,\omega)$. The original data have not been deformed. We have only
put their total charge in its rest frame; the angular function $f_*$
may still change sign and need not be small.

\subsection{Sphere geometry and the constraints}

The collar will change both the leaf geometry and the second tensor.
To see which changes can preserve DEC, let $(G,k)$ be smooth data on a
sphere-foliated portion of the end, with the physical constraint
normalization of the introduction. Write
\begin{equation}\label{eq:foliation}
 \begin{split}
 G&=U^{-2}\,\mathrm dr^2
   +\gamma_{AB}(\mathrm d\omega^A+s^A\mathrm dr)
                 (\mathrm d\omega^B+s^B\mathrm dr),\\
 N&=U(\partial_r-s),\qquad U>0.
 \end{split}
\end{equation}
Here $\gamma$ is the induced sphere metric, $s$ is a sphere-tangential
vector field, and $N$ is the unit normal toward increasing $r$.
Unless otherwise stated, norms, contractions, traces, and differential
operators on a leaf are taken with $\gamma$; in particular,
$\Delta_\gamma=\diver_\gamma\nabla^\gamma$. Denote the leaf second
fundamental form, mean curvature, and normal acceleration by
\begin{equation}\label{eq:leaf-geometry}
 \chi=\frac U2(\partial_r\gamma-\mathcal L_s\gamma),
 \qquad H=\tr_\gamma\chi,
 \qquad B_N=\nabla^G_NN=\nabla^\gamma\log U.
\end{equation}
Decompose $k$ into its normal, mixed, and tangential parts:
\begin{equation}\label{eq:k-decomposition}
 \begin{gathered}
 L=k(N,N),\qquad \eta=k(N,\cdot)|_{T\mathbb S^2},\qquad
 k_T=\frac p2\gamma+\tau,\qquad \tr_\gamma\tau=0,\\
 \widehat\chi=\chi-\frac H2\gamma.
 \end{gathered}
\end{equation}
The following scalar variables will be central:
\begin{equation}\label{eq:mass-variables}
 \begin{aligned}
 d&=\frac r4\tr_\gamma(\partial_r\gamma-\mathcal L_s\gamma),
 &u&=\frac{rH}{2}=Ud,\\
 v&=\frac{rp}{2},
 &F&=\frac r2(1+v^2-u^2).
 \end{aligned}
\end{equation}
Here $u$ is the rescaled mean curvature and $v$ is the rescaled
tangential trace of $k$. The quantity $F$ is their mass combination;
it agrees with the constant mass in the radial model. We use these
variables only where $u,d>0$. The scalar $d$ records the effect of
the nonround leaves on the relation $u=Ud$. It is distinct from the
differential $\mathrm d$; write $\mathrm d_{\mathbb S^2}$ for the
angular differential.

\begin{lemma}[Constraint decomposition]\label{lem:constraints}
For the smooth data $(G,k)$ in \eqref{eq:foliation}, assume $u,d>0$
and use the decomposition \eqref{eq:k-decomposition}. Let
$C=8\pi\mu$, $Q=8\pi J(N)$, and $Z=8\pi J|_{T\mathbb S^2}$.
Then
\begin{equation}\label{eq:constraints}
 \begin{split}
 C={}&\frac{R_\gamma}{2}-NH-U\Delta_\gamma(U^{-1})
      -\frac34H^2-\frac12|\widehat\chi|^2
      +Lp+\frac14p^2-\frac12|\tau|^2-|\eta|^2,\\
 Q={}&H(L-p/2)-\widehat\chi:\tau-Np
      +\diver_\gamma\eta-2\eta(B_N),\\
 Z={}&U(\partial_r\eta-\mathcal L_s\eta)+H\eta
      +\diver_\gamma\tau-\mathrm d_{\mathbb S^2}(L+p/2)
      +(L\gamma-k_T)(B_N,\cdot).
 \end{split}
\end{equation}
In particular, the dominant energy condition is exactly
$C\ge\sqrt{Q^2+|Z|_\gamma^2}$. Moreover, $L$ cancels from the combination
\begin{equation}\label{eq:mass-identity}
 \begin{split}
 C-\frac vuQ={}&\frac{2NF}{ur^2}
  +\left(\frac{R_\gamma}{2}-\frac1{r^2}\right)
  +\left(\frac Uu-1\right)\frac{2-6F/r}{r^2}
  -U\Delta_\gamma(U^{-1})\\
 &-\frac12\bigl(|\widehat\chi|^2+|\tau|^2\bigr)
  +\frac vu\widehat\chi:\tau-|\eta|^2
  -\frac vu\bigl(\diver_\gamma\eta-2\eta(B_N)\bigr).
 \end{split}
\end{equation}
\end{lemma}

\begin{proof}
The first formula of \eqref{eq:leaf-geometry} is the tangential restriction
of $\mathcal L_NG/2$. For a tangential vector field $Y$,
$[N,Y]$ has normal component $-Y(\log U)N$. Metric compatibility and
torsion freeness then give $G(\nabla_NN,Y)=Y(\log U)$, proving the
acceleration formula, including its sign.

The normal Ricci identity and the Gauss equation give, respectively,
\[
 \operatorname{Ric}_G(N,N)=\diver_G B_N-NH-|\chi|^2,
 \qquad
 R_G=R_\gamma+2\operatorname{Ric}_G(N,N)+|\chi|^2-H^2.
\]
Since $B_N$ is tangential,
\[
 \diver_G B_N
 =\diver_\gamma B_N-|B_N|^2
 =-U\Delta_\gamma(U^{-1}).
\]
Thus
\[
 R_G=R_\gamma-2NH-\frac32H^2-|\widehat\chi|^2
                 -2U\Delta_\gamma(U^{-1}).
\]
Also
\[
 (\tr_G k)^2-|k|_G^2
 =2Lp+\frac12p^2-|\tau|^2-2|\eta|^2.
\]
Their half-sum is the asserted expression for $C$.

Here are the momentum calculations, with the normal-acceleration terms
retained. The one-form $k(N,\cdot)$ on the ambient space is
$L N^\flat+\eta$, where $\eta(N)=0$. Its divergence is
\[
 \diver_G\bigl(k(N,\cdot)\bigr)
 =NL+HL+\diver_\gamma\eta-\eta(B_N).
\]
To obtain $(\diver_G k)(N)$, subtract the contraction of $k$ with
$\nabla N$, namely $\eta(B_N)+\chi:k_T$. Subtracting also
$N(\tr_Gk)=N(L+p)$ gives
\[
 Q=HL-\chi:k_T-Np+\diver_\gamma\eta-2\eta(B_N),
\]
which is the second identity in \eqref{eq:constraints}.

For the tangential component, choose a leaf-orthonormal frame
$e_1,e_2$ at the point of calculation and a tangential vector $Y$.
The tangential terms in $(\diver_Gk)(Y)$ are
\[
 \sum_{a=1}^2(\nabla^G_{e_a}k)(e_a,Y)
 = (\diver_\gamma k_T)(Y)+H\eta(Y)
     +\eta\bigl(\chi(Y,\cdot)^\sharp\bigr).
\]
Writing $X=\partial_r-s$, one has
\[
 \nabla^G_NY
 =\chi(Y,\cdot)^\sharp+U[X,Y]-Y(\log U)N.
\]
Consequently, the normal term is
\[
 \begin{split}
 (\nabla^G_Nk)(N,Y)
 ={}&U(\partial_r\eta-\mathcal L_s\eta)(Y)
       -\eta\bigl(\chi(Y,\cdot)^\sharp\bigr)\\
    &+(L\gamma-k_T)(B_N,Y).
 \end{split}
\]
Adding these formulas cancels the two contractions with $\chi$.
Finally subtract $Y(L+p)$ and use
$\diver_\gamma k_T=\tfrac12\mathrm d_{\mathbb S^2}p
+\diver_\gamma\tau$. This proves the formula for $Z$.

To verify \eqref{eq:mass-identity}, first note that $p/H=v/u$, so that
the terms involving $L$ cancel. Set $A=u^2-v^2=1-2F/r$.
The squared trace terms become $-3A/r^2$, whereas $Nr=U$ gives
\[
 -NH+\frac vuNp=-\frac{NA}{ur}+\frac{2AU}{ur^2}.
\]
Substituting $NA=-2NF/r+2FU/r^2$, their sum is
\[
 \frac{2NF}{ur^2}-\frac1{r^2}
       +\left(\frac Uu-1\right)\frac{2-6F/r}{r^2}.
\]
The remaining terms of \eqref{eq:constraints} are precisely those in
\eqref{eq:mass-identity}.
\end{proof}

The identities separate the two adjustments needed in the collar.
Changing $L$ changes the constraints without introducing a radial
derivative of $L$. Changing $F$ produces the first-derivative term
$2NF/(ur^2)$ in \eqref{eq:mass-identity}, where $L$ has canceled.
We will use a radial addition to $L$ for the large null component of
the constraint increment and a slowly increasing addition to $F$ for
the small component. Their sizes will also control the angular
momentum increment.

\subsection{The original data in foliation variables}

Before changing the data, we identify the sizes of these variables.
Put a subscript $0$ on the quantities associated with $G_0=g$ and
$k_0=K$. In particular, the limit of $F_0$ must contain the leading
coefficients of both tensors, because it is the Bondi mass aspect
that must survive the replacement.

\begin{lemma}\label{lem:original-asymptotics}
For smooth data in the class \eqref{eq:cks}, expressed in the balanced
chart of Lemma~\ref{lem:balance}, the original foliation variables satisfy
\begin{equation}\label{eq:original-asymptotics}
 \begin{aligned}
 r^{-2}\gamma_0
  &=\sigma+\frac{m^g}{r^3}+O_3(r^{-4}),
       &s_0&=O_3(r^{-5}),\\
 d_0&=1-\frac{3\tr_\sigma m^g}{4r^3}+O_2(r^{-4}),\\
 \frac{U_0}{\sqrt{1+r^2}}
  &=1-\frac{m^r}{2r^3}+O_2(r^{-4}),\\
 \frac{u_0}{\sqrt{1+r^2}}
  &=1-\frac{m^r/2+(3/4)\tr_\sigma m^g}{r^3}
       +O_2(r^{-4}),\\
 \frac{v_0}{r}
  &=1+\frac{\tr_\sigma(m^K-m^g)}{2r^3}+O_2(r^{-4}),
       &F_0&=f_*+O_2(r^{-1}),\\
 L_0-1&=O_2(r^{-3}),
       &r^{-1}\eta_0&=O_2(r^{-3}),\qquad
         r^{-2}\tau_0=O_2(r^{-3}).
 \end{aligned}
\end{equation}
In particular, $d_0,u_0,v_0>0$ sufficiently far out.
\end{lemma}

\begin{proof}
The tangential metric is $\gamma_0=g_{AB}$, and
$s_0^A=\gamma_0^{AB}g_{rB}$, which proves the first line. Writing
$q_0=r^{-2}\gamma_0$, the definition of $d_0$ becomes
\[
 d_0=1+\frac r4\tr(q_0^{-1}\partial_rq_0)
                  -\frac r2\diver_{\gamma_0}s_0.
\]
The first correction has leading coefficient
$-3\tr_\sigma m^g/(4r^3)$; the shift term is $O_2(r^{-4})$.
This calculation consumes one of the three allowed metric derivatives
and leaves the stated two differentiated orders.

The normal lapse follows from
\[
 U_0^{-2}=g_{rr}-\gamma_0(s_0,s_0)
 =\frac1{1+r^2}+\frac{m^r}{r^5}+O_3(r^{-6}),
\]
where $\gamma_0(s_0,s_0)=O_3(r^{-8})$. Expanding the inverse square root
gives the displayed formula for $U_0$; multiplication by $d_0$ gives
the formula for $u_0$.

Next, taking the tangential trace of $K$ gives
\[
 p_0=\tr_{\gamma_0}K_T
 =2+\frac{\tr_\sigma(m^K-m^g)}{r^3}+O_2(r^{-4}),
\]
which proves the expansion of $v_0$. If
$\alpha=m^r/2+(3/4)\tr_\sigma m^g$ and
$\beta=\tfrac12\tr_\sigma(m^K-m^g)$, these formulas imply
\[
 u_0^2=1+r^2-\frac{2\alpha}{r}+O_2(r^{-2}),\qquad
 v_0^2=r^2+\frac{2\beta}{r}+O_2(r^{-2}).
\]
Since $\alpha+\beta=M/4$, the definition of $F_0$ yields
$F_0=M/4+O_2(r^{-1})$, with the same two differentiated orders.

Finally, use $K-g=b-a$ in the normal, mixed, and trace-free tangential
components of \eqref{eq:k-decomposition}. The normal component is
$O_2(r^{-3})$, the mixed angular covector is $O_2(r^{-2})$, and the
trace-free angular tensor is $O_2(r^{-1})$. The corresponding components
of $g$ are $1$, $0$, and $0$. This gives the last line of
\eqref{eq:original-asymptotics}; no leading coefficient for $b_{rr}$ is
required.
\end{proof}
